\documentclass[11pt]{article}
\usepackage[margin=1in]{geometry}
\usepackage[T1]{fontenc}
\usepackage[utf8]{inputenc}
\usepackage{lmodern}
\usepackage{amsmath,amssymb}
\usepackage{booktabs}
\usepackage{array}
\usepackage{enumitem}
\usepackage{hyperref}
\usepackage{xcolor}

\hypersetup{colorlinks=true,linkcolor=blue,citecolor=blue,urlcolor=blue}
\setlist{nosep}

\title{MATVERSE 2.0: A Constitutional Architecture for Governed Informational Transformation in Federated Research Systems}
\author{Mateus Alves Ar\^eas\\ORCID: 0009-0008-2973-4047}
\date{Manuscript candidate v0.1 --- September 2026}

\begin{document}
\maketitle

\begin{abstract}
Scientific AI systems increasingly combine interpretation, decision support, tool execution, memory, and automated knowledge transformation. When these functions are concentrated within the same operational authority, it becomes difficult to distinguish what was proposed, authorized, executed, preserved, or subsequently reconstructed. We introduce MATVERSE 2.0, a constitutional architecture for governed informational transformations in federated research systems. The architecture explicitly separates identity, authorship, semantic transformation, admissibility, execution, persistent memory, provenance, and replay. It further represents epistemic status and coherence as multidimensional structures rather than single confidence scores. We describe the constitutional model, its federated authority structure, the $\Omega$-Gate admissibility mechanism, EvidenceOS provenance and replay, and a reference end-to-end implementation that distinguishes decision receipt, execution receipt, memory commitment, and replay receipt. A local vertical slice demonstrates fail-closed behavior under missing authorship and exact reconstruction within the declared deterministic scope. The work makes no claim that replay establishes ontological informational primacy or that the architecture constitutes biological life, consciousness, clinical efficacy, or quantum advantage. Instead, MATVERSE is presented as an engineering and research infrastructure for studying how informational transformations can remain governable, attributable, and reconstructible as scientific systems increase in complexity.
\end{abstract}

\section{Introduction}

Scientific systems increasingly transform information before humans can independently inspect every intermediate state, making the governance of transformation itself a scientific and engineering problem. The central challenge is therefore not only whether a system can generate useful results, but whether the path from information to action remains attributable, admissible, inspectable, and reconstructible.

Modern AI-assisted research systems can combine retrieval, interpretation, synthesis, decision support, tool invocation, persistence, and automated reporting in one operational loop. This concentration of capability creates a structural ambiguity. A generated statement can be mistaken for an authorized decision; an authorized decision can be mistaken for an executed action; an executed action can be mistaken for a persistent scientific result; and a persisted record can later be rewritten or summarized in a form that obscures its origin. The resulting failure is not necessarily a model-quality failure. It is often an authority, provenance, and state-separation failure.

MATVERSE 2.0 addresses this problem through a constitutional architecture that separates interpretation, authorization, execution, persistent memory, and replay. The architecture treats these operations as distinct state transitions governed by explicit identities, policies, provenance records, and receipts. Its purpose is not to determine universal truth. Its purpose is to make transformations of scientific information governable and reconstructible within a declared computational scope.

The contribution of this paper is intentionally narrower than the historical MATVERSE research program. We present MATVERSE as (i) an architecture, (ii) a method for governed informational transformation, (iii) a reference implementation, and (iv) an extensible research program. We do not require acceptance of broader hypotheses concerning digital life, consciousness, fundamental informational ontology, clinical efficacy, or quantum advantage.

The paper makes five main contributions:

\begin{enumerate}
  \item \textbf{Constitutional separation.} A system architecture in which interpretation, admissibility, execution, memory commitment, and replay are represented as distinct operations rather than a single agent action.
  \item \textbf{Federated epistemic authority.} A model that separates authorship, source class, domain authority, authorization, AI synthesis, and shared human--system decisions.
  \item \textbf{Multidimensional coherence.} A representation in which semantic, logical, evidential, causal, architectural, operational, external, and other dimensions can retain independent states instead of being collapsed into a single truth score.
  \item \textbf{Evidence and replay architecture.} Explicit separation between decision receipts, execution receipts, memory commits, and replay receipts, preserving the distinction between admission, action, persistence, and reconstruction.
  \item \textbf{Reference vertical slice.} A local end-to-end implementation in which a candidate object traverses identity/provenance, semantic qualification, memory structuring, adversarial review, policy gating, execution, commitment, evidence recording, and replay.
\end{enumerate}

The central claim is correspondingly bounded: \emph{governed informational transformations in federated research systems become more auditable when interpretation, admissibility, execution, memory, and replay are explicit, separately attributable operations with preserved provenance and policy state.}

\section{Problem Statement and Scope}

\subsection{Concentration of epistemic and operational authority}

Consider a system that can generate an interpretation, judge that interpretation, authorize an action, execute it, persist a record, and later rewrite the record. Even if every individual component is reliable, the combined system lacks a clear answer to several questions:

\begin{itemize}
  \item Who authored the input and who transformed it?
  \item Which policy authorized the action?
  \item Was a decision merely emitted, or was an external effect executed?
  \item Was the result committed to persistent memory?
  \item Can the recorded state reconstruct the declared execution?
  \item Did later summarization alter the evidential meaning of the original object?
\end{itemize}

MATVERSE treats these questions as architectural invariants rather than optional metadata. The design objective is to prevent one stage from silently acquiring the authority of another.

\subsection{Operational versus fundamental programs}

MATVERSE historically contains both operational and foundational research questions. This paper separates them.

\paragraph{Operational program.} The operational objective is to build and study transformations that preserve identity, context, authorship, applicable constraints, uncertainty, provenance, and reconstructibility.

\paragraph{Fundamental program.} Questions concerning the ontological or causal role of information in physical reality remain open research hypotheses. They are not premises of the architecture and are not validated by the experiments reported here.

This distinction is essential. Computational replay can provide engineering evidence of reconstructibility or deterministic equivalence within a model; it does not by itself establish ontological precedence of information over physical processes.

\section{Related Work and Claim Boundary}

MATVERSE intersects several established areas, and its claims must be stated relative to them.

\subsection{Provenance models}

The W3C PROV family defines a domain-agnostic model for representing entities, activities, agents, derivations, and responsibility in provenance records~\cite{w3cprov}. MATVERSE does not claim invention of provenance. Instead, it treats provenance as one constitutional dimension in a runtime that also separates authorization, execution, memory commitment, and replay.

\subsection{Research objects and reproducible publications}

RO-Crate provides a practical JSON-LD packaging model for research objects, contextual entities, files, software, workflows, authors, and licenses~\cite{rocrate}. Whole Tale connects publications to data, code, computational environments, and workflows, explicitly pursuing reproducible and ``living'' computational publications~\cite{wholetale}. These systems establish strong prior art for packaging research artifacts and linking publications to executable context. MATVERSE therefore does not claim novelty for the general concept of a research object or living publication.

\subsection{Software supply-chain attestations}

SLSA and related in-toto attestation mechanisms formalize provenance for software artifacts, including inputs, builders, build definitions, and verifiable properties~\cite{slsa}. This demonstrates that cryptographically anchored provenance and attestations are established engineering mechanisms. MATVERSE uses similar integrity concepts but applies them to governed scientific transformation and explicit epistemic/authority states rather than only software build provenance.

\subsection{Governance-first agent architectures}

Recent work makes generic ``separation of powers'' claims insufficient as a novelty basis. In particular, LATTICE presents a governance-first architecture for autonomous AI operations that separates planning, execution, and governance, implements policy-as-code enforcement, and produces cryptographic audit trails~\cite{lattice}. MATVERSE therefore does not claim that separating governance from execution is itself novel. The narrower research question pursued here is how such separation composes with scientific-object identity, authorship/source classes, multidimensional epistemic state, distinct receipts for decision/execution/commit/replay, and reconstruction of governed transformations.

This related-work boundary is deliberate: each individual mechanism has substantial precedent. The object of study is the composition and operational semantics used to govern transformations in a federated scientific system.

\section{Design Principles}

MATVERSE 2.0 is defined by the following design principles.

\begin{enumerate}
  \item \textbf{Identity preservation.} Objects and components maintain explicit identities across transformations.
  \item \textbf{Authorship preservation.} A linguistic or computational transformation does not silently change authorship.
  \item \textbf{Domain-local authority.} Expertise or authority in one domain does not imply universal authority.
  \item \textbf{Separation of powers.} Generation, admissibility, execution, persistence, and replay are distinct capabilities.
  \item \textbf{Fail-closed admissibility.} Missing required evidence, identity, authorship, or policy state must not be silently treated as permission.
  \item \textbf{Uncertainty preservation.} Unknown or unresolved states remain explicit rather than being coerced into pass/fail certainty.
  \item \textbf{Provenance by construction.} Transformations emit records sufficient to identify inputs, policy, decision, and downstream effects within the declared scope.
  \item \textbf{Replay as a separate state.} A successful decision or execution is not equivalent to successful reconstruction.
  \item \textbf{No silent canonical promotion.} Repetition, summarization, or AI synthesis does not automatically promote a claim to stronger epistemic status.
\end{enumerate}

These principles are constitutional in the software-architecture sense: components can evolve, but the authority boundaries and invariants constrain valid state transitions.

\section{Constitutional Architecture}

The reference architecture is organized as a federated chain rather than a monolithic agent. The minimal path used in the reference vertical slice is:

\begin{verbatim}
SOURCE
  -> ATLAS / REGISTRY
  -> CASSANDRA
  -> MNB-CANDIDATE
  -> COHERENCE / REDTEAM
  -> OMEGA-GATE
  -> DECISION RECEIPT
  -> EXECUTOR
  -> EXECUTION RECEIPT
  -> MEMORY COMMIT
  -> EVIDENCEOS
  -> REPLAY
  -> REPLAY RECEIPT
\end{verbatim}

\subsection{Atlas / Registry}

Atlas represents identity, relationships, lineage, conflicts, and registry-level metadata. Its constitutional role is not to decide truth but to preserve where an object is situated and what identities or dependencies are relevant to later decisions.

\subsection{Cassandra}

Cassandra performs semantic qualification, normalization, and source/claim structuring. It can classify an input as, for example, \texttt{AI\_SYNTHESIS} or another source class. Its output is a qualified object, not a truth verdict.

\subsection{MNB candidate}

The MNB layer represents a candidate memory object with contextualized content, provenance, state, and integrity information. A candidate is not committed merely because it exists. Promotion to persistent governed memory is a later operation.

\subsection{Coherence and REDTEAM}

This stage performs consistency and adversarial checks before admission. Importantly, coherence is not defined as universal truth. It is a structured evaluation over declared dimensions and invariants.

\subsection{$\Omega$-Gate}

The $\Omega$-Gate evaluates admissibility under a versioned policy and emits a localized decision such as \texttt{PASS}, \texttt{HOLD}, \texttt{BLOCK}, or \texttt{ESCALATE}. The gate decides whether a specific transition is admissible under a specific policy; it does not establish that an arbitrary proposition is universally true.

\subsection{Executor}

The executor performs only actions that have crossed the relevant authority boundary. This preserves the distinction between a decision and an external or internal effect.

\subsection{EvidenceOS and replay}

EvidenceOS preserves the event chain and receipts required to inspect and reconstruct the declared execution. Replay is a later verification/reconstruction operation. It may produce exact agreement, divergence, or another explicit comparison result.

\section{Informational Transformation Model}

We model a governed transformation as

\begin{equation}
T:(S_t,C_t,A_t,P_t)\rightarrow(S_{t+1},R_t),
\end{equation}

where $S_t$ is informational state, $C_t$ is context, $A_t$ represents authorship and applicable authority, $P_t$ is the versioned policy, and $R_t$ is the resulting receipt or transformation record.

A transformation is not admissible merely because a candidate next state can be computed. Admissibility requires satisfaction of the constitutional conditions relevant to the transition. We can represent this generically as

\begin{equation}
\mathrm{Admit}(T)=1 \quad \Leftrightarrow \quad \bigwedge_{i\in I_T} I_i(S_t,C_t,A_t,P_t)=1,
\end{equation}

where $I_T$ is the set of invariants applicable to transformation $T$. The formulation is intentionally policy-relative: weights, thresholds, and required invariants are versioned configuration unless independently established as domain laws.

The broader informational cycle used by MATVERSE is an architectural model:

\begin{verbatim}
INFORMATION -> REPRESENTATION -> TRANSFORMATION -> ADMISSIBILITY
            -> EXECUTION -> OBSERVATION -> EVIDENCE -> REPLAY
            -> NEW INFORMATION
\end{verbatim}

This sequence is not presented as a universal law of nature. It is a system-design decomposition for tracing scientific transformations.

\section{Epistemic and Authorship Model}

MATVERSE separates source/epistemic roles that are often conflated in AI-assisted research:

\begin{center}
\begin{tabular}{p{0.25\linewidth}p{0.62\linewidth}}
\toprule
\textbf{Class} & \textbf{Meaning} \\
\midrule
\texttt{HUMAN\_SOURCE} & Content directly attributable to a human source. \\
\texttt{AI\_SYNTHESIS} & A model-generated synthesis, transformation, or interpretation. \\
\texttt{SHARED\_DECISION} & A decision explicitly produced through a defined human--system authority process. \\
\texttt{EXPERIMENTAL\_RESULT} & An output tied to an instrument, protocol, execution, and evidence record. \\
\texttt{EXTERNAL\_SCIENCE} & A result or claim attributable to an external scientific source. \\
\bottomrule
\end{tabular}
\end{center}

Three rules follow.

First, repetition of a claim does not create independent evidence. Second, transformation of language does not automatically change authorship. Third, architectural authority does not imply scientific authority over every domain. These rules are designed to preserve uncertainty and provenance across heterogeneous contributors and agents.

\section{Multidimensional Coherence and Admissibility}

A single scalar ``truth score'' is generally insufficient to represent the state of a scientific object. MATVERSE instead permits a vector of partially independent dimensions, for example

\begin{equation}
\Psi(x)=\langle \psi_{sem},\psi_{log},\psi_{evid},\psi_{causal},\psi_{arch},\psi_{auth},\psi_{temp},\psi_{op},\psi_{ext},\psi_{human},\psi_{unc}\rangle.
\end{equation}

The elements need not share one numeric scale. A practical object may simultaneously be semantically coherent, architecturally valid, operationally demonstrated locally, evidentially incomplete, and externally unreplicated. For example:

\begin{verbatim}
semantic:      PASS
logical:       PASS
evidential:    HOLD
causal:        UNKNOWN
architectural: PASS
operational:   PASS_LOCAL
external:      UNKNOWN
\end{verbatim}

This representation prevents internal coherence from being silently promoted to external scientific verification. The $\Omega$-Gate consumes the dimensions relevant to a particular transition under a versioned policy; it need not collapse the entire vector into a universal score.

\section{Receipt Separation and Replay}

A core invariant is

\begin{equation}
\mathrm{PASS}\neq\mathrm{EXECUTED}\neq\mathrm{COMMITTED}\neq\mathrm{REPLAYED}.
\end{equation}

We distinguish four records:

\begin{align}
R_d &= \text{decision receipt},\\
R_e &= \text{execution receipt},\\
M_c &= \text{memory commit},\\
R_r &= \text{replay receipt}.
\end{align}

A simplified causal sequence is

\begin{equation}
S_t \xrightarrow{G} R_d \xrightarrow{E} R_e \xrightarrow{C} M_c \xrightarrow{P} R_r,
\end{equation}

where $G$ is the gate, $E$ is execution, $C$ is commit, and $P$ is replay. Each transition has a distinct evidential meaning.

A decision receipt records that a policy emitted a decision. It does not prove that execution occurred. An execution receipt records the execution event within the defined runtime scope. A memory commit records persistence. A replay receipt records the outcome of reconstruction/comparison. This separation is particularly important for external side effects, where uncertain execution state must not be treated as safe to retry automatically.

\section{Reference Implementation}

The reference implementation is intentionally treated as an engineering artifact rather than proof of the entire research program. The executed local vertical slice traversed the following sequence:

\begin{verbatim}
ATLAS -> CASSANDRA -> MNB-CANDIDATE -> COHERENCE/REDTEAM
      -> OMEGA-GATE -> DECISION RECEIPT -> EXECUTOR
      -> EXECUTION RECEIPT -> MEMORY COMMIT -> EVIDENCEOS
      -> REPLAY
\end{verbatim}

The recorded local run classified the source as \texttt{AI\_SYNTHESIS}, created and hashed an MNB candidate, completed six of six declared coherence/REDTEAM checks, passed the $\Omega$-Gate under policy \texttt{MV-POL-ARCH-INTEGRATION-1.0.0}, emitted decision and execution receipts, persisted a memory commit, recorded four chained EvidenceOS events, and reported \texttt{EXACT\_MATCH} across ten of ten compared replay stages.

These observations are explicitly local. They establish that the declared components can participate in an end-to-end governed execution within the tested environment. They do not establish independent third-party reproduction or universal properties of the architecture.

\section{Evaluation}

The Paper 1 evaluation is designed around constitutional properties rather than aggregate model performance.

\subsection{E1: Valid-path execution}

For a valid input with identity, authorship, source, and policy requirements satisfied, the expected sequence is

\begin{verbatim}
PASS -> EXECUTION -> COMMIT -> REPLAY.
\end{verbatim}

\textbf{Observed local result:} the reference vertical slice traversed the complete chain, emitted distinct receipts, committed memory, and produced an exact replay comparison within the declared scope. Status: \texttt{PASS\_LOCAL}.

\subsection{E2: Missing authorship}

The adversarial perturbation removes required authorship from the candidate object.

\textbf{Expected:} \texttt{BLOCK} before execution.

\textbf{Observed local result:} the gate returned \texttt{BLOCK}, emitted a decision record, and prevented downstream execution. This demonstrates fail-closed behavior for this tested invariant. It does not establish fail-closed behavior for every possible missing field or adversarial condition.

\subsection{E3: Tampered artifact}

The planned perturbation changes an input or persisted artifact between execution and replay.

\textbf{Expected:} divergence or integrity failure must be surfaced rather than silently accepted.

\textbf{Current status:} specified in the Paper 1 test plan; not promoted here as an observed result unless accompanied by the corresponding frozen evidence artifact.

\subsection{E4: Policy mismatch}

The planned perturbation changes the policy version or policy fingerprint between decision and later evaluation.

\textbf{Expected:} \texttt{BLOCK} or \texttt{HOLD} according to the frozen contract.

\textbf{Current status:} specified; not claimed as an observed Paper 1 result in this manuscript candidate.

\subsection{E5: Independent replay}

The current implementation supports replay outside the original transformation step and records exact local comparison for the reference slice.

\textbf{Current status:} local independent-process reconstruction is within the demonstrated program; independent third-party reproduction remains \texttt{HOLD}.

\begin{table}[h]
\centering
\small
\begin{tabular}{p{0.16\linewidth}p{0.27\linewidth}p{0.25\linewidth}p{0.18\linewidth}}
\toprule
\textbf{Test} & \textbf{Perturbation / condition} & \textbf{Observed state} & \textbf{Status} \\
\midrule
E1 & Valid identity, authorship, policy, source & End-to-end traversal; exact local replay & PASS\_LOCAL \\
E2 & Missing authorship & BLOCK before execution & PASS\_LOCAL \\
E3 & Artifact tampering & Not promoted as observed result & SPECIFIED \\
E4 & Policy mismatch & Not promoted as observed result & SPECIFIED \\
E5 & Replay outside original step & Local exact comparison; no third-party replication & PASS\_LOCAL / HOLD\_EXTERNAL \\
\bottomrule
\end{tabular}
\caption{Paper 1 constitutional evaluation states. ``Specified'' is deliberately distinct from ``observed.''}
\end{table}

\section{Discussion}

\subsection{What the reference slice demonstrates}

The reference run supports a narrow but useful conclusion: the architecture can implement a governed path in which source qualification, admissibility, execution, persistence, and replay are distinct events with distinct records. The missing-authorship perturbation further demonstrates that at least one required authority invariant can fail closed before execution.

The exact replay result is engineering evidence of reconstruction within the tested deterministic scope. It is not evidence that all relevant real-world phenomena are determined by the recorded information, nor that information has ontological priority over physical processes.

\subsection{Why multidimensional state matters}

Scientific systems routinely encounter objects that are internally coherent but externally unverified, executable but scientifically underdetermined, or reproducible locally but not independently replicated. A scalar confidence score tends to obscure these differences. MATVERSE instead makes such partial states explicit and permits policy to depend only on the dimensions relevant to the attempted transition.

\subsection{Federation rather than monolithic agency}

MATVERSE uses federation to localize authority. A component that can synthesize text is not automatically allowed to execute tools; a component that can execute does not automatically decide admissibility; a registry that tracks lineage does not become a scientific judge. This resembles separation-of-concerns patterns in provenance, supply-chain security, and governance-first agent architectures, but is applied here to the lifecycle of scientific informational transformations.

\subsection{Downstream governed effectors}

The same authority pattern can be applied to external effectors such as a publication runtime. A publication agent can organize metadata and prepare a frozen package without having authority to publish it; a human approval can delegate authority only for an exact package; and external identity can later be reconciled into the scientific object's state. Such applications are downstream demonstrations of the constitutional pattern and are not treated as evidence for stronger organism-level claims.

\section{Limitations}

This work has several explicit limitations.

\begin{enumerate}
  \item \textbf{Local evidence boundary.} The reference vertical slice is a local implementation result. Independent third-party reproduction is not yet claimed.
  \item \textbf{Partial adversarial coverage.} Missing authorship has an observed fail-closed result, while other planned perturbations require separately frozen evidence before being reported as results.
  \item \textbf{No universal metric claim.} Coherence dimensions and policy thresholds are operational constructs unless independently calibrated. No universal truth score is claimed.
  \item \textbf{No ontological claim.} Replay and provenance do not establish that information is ontologically prior to matter.
  \item \textbf{No biological-life or consciousness claim.} Biological metaphors used elsewhere in the research program are not results of this paper.
  \item \textbf{No clinical claim.} Human-system research programs such as NeuroFlow are outside the demonstrated scope.
  \item \textbf{No quantum-advantage claim.} Quantum-related research programs are not evidence for this architectural contribution.
  \item \textbf{Novelty requires decomposition.} Provenance, research objects, living publications, separation of governance from execution, cryptographic attestations, and replay-related ideas all have prior art. The defensible contribution must be assessed at the level of the specific composition and semantics implemented here.
\end{enumerate}

\section{Research Programs Beyond Paper 1}

Several MATVERSE directions remain active but are not dependencies of the founder paper. These include formalization of additional memory units and causal inheritance, organism-level hypotheses, human--machine coadaptation, NeuroFlow, quantum-classical workbenches, and fundamental informational models. Their inclusion in the broader program does not promote them to demonstrated results.

The architecture is designed to preserve this distinction: a research hypothesis can remain registered, versioned, and traceable without being admitted as established fact.

\section{Reproducibility and Artifact Availability}

The intended public release couples four versioned objects:

\begin{enumerate}
  \item the MATVERSE founder paper;
  \item the MATVERSE public canon;
  \item the reference implementation; and
  \item an evidence pack containing protocols, inputs, tests, outputs, manifests, hashes, receipts, and replay material.
\end{enumerate}

The current development repository is available at
\href{https://github.com/MatVerse-py/Gpt-project-bridge}{github.com/MatVerse-py/Gpt-project-bridge}.
A release tag, archival identifier, and immutable evidence-pack references must be frozen before final submission. No DOI or independent replication status is asserted in this manuscript candidate.

\section{Conclusion}

MATVERSE 2.0 presents a constitutional architecture for governed informational transformation in federated research systems. Its central engineering choice is to prevent interpretation, admissibility, execution, persistent memory, and replay from collapsing into one authority or one undifferentiated state. The architecture preserves authorship and source class, maintains multidimensional epistemic state, uses versioned policy gates, distinguishes decision/execution/commit/replay receipts, and supports reconstruction of a declared execution path.

A local reference vertical slice demonstrates that these separations can be implemented end-to-end within the tested scope, including fail-closed rejection when authorship is missing and exact local replay across the recorded comparison stages. The result is not a universal truth engine and not evidence of digital life or fundamental informational ontology. It is a concrete systems contribution: an auditable framework for governing how scientific information is transformed into decisions, actions, persistent state, and reproducible records.

\begin{thebibliography}{9}

\bibitem{w3cprov}
L. Moreau and P. Missier (eds.).
\newblock \emph{PROV-DM: The PROV Data Model}.
\newblock W3C Recommendation, 30 April 2013.
\newblock \url{https://www.w3.org/TR/prov-dm/}.

\bibitem{rocrate}
P. Sefton, E. \`O Carrag\'ain, S. Soiland-Reyes, et al.
\newblock \emph{RO-Crate Metadata Specification}.
\newblock Research Object Crate community, current specification family.
\newblock \url{https://www.researchobject.org/ro-crate/specification.html}.

\bibitem{wholetale}
B. Ludaescher, K. Chard, N. Gaffney, M. B. Jones, J. Nabrzyski, V. Stodden, and M. Turk.
\newblock Capturing the ``Whole Tale'' of Computational Research: Reproducibility in Computing Environments.
\newblock \emph{arXiv:1610.09958}, 2016; related environment architecture published in \emph{Future Generation Computer Systems}, 94:854--867, 2019.

\bibitem{slsa}
SLSA Framework.
\newblock \emph{Supply-chain Levels for Software Artifacts, Specification v1.2}.
\newblock \url{https://slsa.dev/spec/v1.2/}.

\bibitem{lattice}
E. Calboreanu.
\newblock LATTICE: a governance-first architecture for authorized autonomous AI operations.
\newblock \emph{Frontiers in Artificial Intelligence}, 9, 2026.
\newblock doi:10.3389/frai.2026.1800407.

\end{thebibliography}

\end{document}
